% MACRO I (ECO387C), Fall 2026 - midterm study packet. Author: Philip Livdan.
% Walks through the class material in lecture order (Lectures 1-4, Discussions
% 1-6, PSets 1-6): every model is set up from scratch and every result derived
% in numbered steps with all algebra displayed, then worked problems follow.
% Every number is checked in practice-packet-checks.py (output in
% practice-packet-checks.txt). Rerun it before changing any number.
% Font gotchas (offline tectonic cache): 12pt only, no math in \section*,
% \mathbf not \boldsymbol, no \widehat in exponents, no monospace.
\documentclass[12pt]{article}
\usepackage{amsmath,amssymb}
\usepackage[margin=1in]{geometry}
\DeclareMathSizes{12}{12}{8}{7}
\DeclareMathSizes{11}{11}{8}{7}
\DeclareMathSizes{10}{10}{8}{7}
\renewcommand{\ttdefault}{lmr}
\setlength{\parindent}{0pt}
\setlength{\parskip}{0.5\baselineskip}
\newcommand{\E}{E}
\DeclareMathOperator{\Var}{Var}
\DeclareMathOperator{\Cov}{Cov}
\newcommand{\st}[1]{\par\smallskip\textbf{Step #1.}\ }
\newcommand{\given}[1]{\textit{Given.} #1\par}
\newcommand{\uses}[1]{\textit{Uses.} #1\par}
\newcommand{\goal}[1]{\textit{Goal.} #1\par}
\newcommand{\ans}[1]{\[\boxed{\ #1\ }\]}
\newcommand{\prob}[1]{\par\bigskip\textbf{Worked problem #1.}\ }
\newcommand{\sol}{\par\smallskip\textbf{Solution.}\par}
\newcommand{\src}[1]{\par{\small\textit{Source in class.} #1}\par}

\title{Midterm Study Packet\\[1mm] \large Macroeconomics I (ECO387C), Fall 2026}
\author{Philip Livdan}
\date{October 6, 2026}

\begin{document}
\maketitle

This packet rebuilds the course from the ground up in the order it was taught.
Each part has three layers. \textbf{Setup} writes down the environment exactly
as in lecture (who, what is chosen, what is given, the constraints).
\textbf{Derivation} gets every result in numbered steps, each saying what is
done and why, with every line of algebra displayed. \textbf{Worked problems}
apply the derivation to exam-style questions. Cover the solution, attempt the
problem, then check line by line. The exam is closed book, so the boxed results
are the things to know cold.

Parts: 0 toolkit, 1 what macro is (Lecture 1), 2 national accounts
(Discussion 1), 3 the endowment economy (Lecture 2, Discussion 2, PSet 2),
4 infinite horizon and uncertainty (Lecture 2, PSets 3 and 6),
5 log-linearization and time averaging (Discussion 3, PSet 3),
6 dynamic programming (Lecture 3, Discussion 4, PSets 3--4),
7 numerical dynamic programming (Discussions 5--6, PSets 4--5),
8 consumption (Lecture 4).

\tableofcontents
\newpage

% ======================================================================
\section*{Part 0. Toolkit}
\addcontentsline{toc}{section}{Part 0. Toolkit}

These five tools are used over and over. Know them before anything else.

\textbf{T1. Lagrangian recipe with inequality constraints (KKT).}
To maximize $f(x)$ subject to $g_i(x)\ge0$:
\begin{enumerate}
\item Write $\mathcal L=f(x)+\sum_i\lambda_i g_i(x)$, one multiplier per constraint.
\item First-order conditions $\partial\mathcal L/\partial x_j=0$ for every choice variable.
\item Constraints hold, $g_i(x)\ge0$.
\item Complementary slackness, $\lambda_i g_i(x)=0$. Either the constraint binds or its multiplier is zero.
\item Non-negativity, $\lambda_i\ge0$.
\end{enumerate}
The multiplier is the marginal value of relaxing the constraint. For a budget
constraint its unit is utils per unit of goods.

\textbf{T2. Geometric sums.} For $|x|<1$,
\[
\sum_{j=0}^{\infty}x^j=\frac{1}{1-x},\qquad
\sum_{j=0}^{T}x^j=\frac{1-x^{T+1}}{1-x},\qquad
\sum_{\substack{t=0\\t\text{ even}}}^{T}x^{t}=\sum_{s=0}^{(T-1)/2}(x^2)^s=\frac{1-x^{T+1}}{1-x^2}\ \ (T\text{ odd}).
\]
The most used special case is
\[
\sum_{j=0}^{\infty}\frac{1}{(1+r)^j}=\frac{1}{1-\frac{1}{1+r}}=\frac{1+r}{(1+r)-1}=\frac{1+r}{r}.
\]

\textbf{T3. Envelope theorem.} If $V(x)=\max_{y}G(x,y)$ with maximizer $y=g(x)$, then
\[
V'(x)=\frac{d}{dx}G\big(x,g(x)\big)=G_x\big(x,g(x)\big)+\underbrace{G_y\big(x,g(x)\big)}_{=0\text{ by the FOC}}g'(x)=G_x\big(x,g(x)\big).
\]
Only the direct effect of the state survives. The indirect effect through the
choice is zero because the choice was already optimal.

\textbf{T4. Taylor expansions.} First order around $\bar x$:
$f(x)\approx f(\bar x)+f'(\bar x)(x-\bar x)$. Second order adds
$\tfrac12f''(\bar x)(x-\bar x)^2$. Also $\ln(1+z)\approx z$ for small $z$, so
$\ln x-\ln\bar x=\ln\big(1+\tfrac{x-\bar x}{\bar x}\big)\approx\tfrac{x-\bar x}{\bar x}$.

\textbf{T5. CRRA algebra.} Course notation (Lecture 4) is
$u(c)=\dfrac{c^{1-1/\sigma}}{1-1/\sigma}$ with $\sigma$ the intertemporal
elasticity of substitution (IES). Psets write $u(c)=\dfrac{c^{1-\gamma}}{1-\gamma}$,
so $\gamma=1/\sigma$ is relative risk aversion. Derivatives:
\[
u'(c)=c^{-\gamma},\qquad u''(c)=-\gamma c^{-\gamma-1},\qquad u'''(c)=\gamma(\gamma+1)c^{-\gamma-2}.
\]
Relative risk aversion and relative prudence are
\[
-\frac{u''(c)c}{u'(c)}=\frac{\gamma c^{-\gamma-1}c}{c^{-\gamma}}=\gamma,\qquad
-\frac{u'''(c)c}{u''(c)}=\frac{\gamma(\gamma+1)c^{-\gamma-2}c}{\gamma c^{-\gamma-1}}=\gamma+1 .
\]
Log utility is the case $\gamma=\sigma=1$.

% ======================================================================
\section*{Part 1. What macroeconomics is (Lecture 1)}
\addcontentsline{toc}{section}{Part 1. What macroeconomics is}

\textbf{Four defining features.}
\begin{enumerate}
\item \emph{Aggregates.} $C_t=\sum_ic_{it}$, $Y_t=\sum_iy_{it}$. We model individuals and add them up.
\item \emph{General equilibrium.} Prices are determined inside the model by market clearing, not taken from outside.
\item \emph{Dynamics.} Decisions today depend on the future (saving, investment), so time is explicit.
\item \emph{Quantitative.} Models are calibrated or estimated and asked for numbers, not just signs.
\end{enumerate}
Core fields are business cycles (with monetary and fiscal policy) and growth.
Applied fields include macro labor, structural change, inequality, and climate.
Adjacent fields are international trade and finance, macro-finance, and
macro-development.

\textbf{History in one line per era.}
\begin{itemize}
\item Classical economics, then the Great Depression and Keynes (aggregate demand). National accounts are built (Part 2).
\item 1940s--60s Keynesian consensus. Hicks's IS-LM in $(r,Y)$, Okun's law (GDP down 2\% goes with unemployment up 1 point), the Phillips curve as a ``menu'' of inflation--unemployment pairs.
\item 1960s--70s monetarism. Friedman and Phelps argue the long-run Phillips curve is vertical. The 1970s oil shocks bring high inflation and high unemployment together, which the menu cannot explain.
\item 1970s--90s new classical economics. Lucas critique and rational expectations, then Kydland and Prescott (1982) and real business cycle (RBC) models. Saltwater versus freshwater schools.
\item 1990s to today. New Keynesian models (Gal\'i, Woodford) add sticky prices to the RBC core, so monetary policy has real effects. Empirical side uses VARs (Sims, Sargent), $\mathbf Y_t=\alpha+B\mathbf Y_{t-1}+u_t$, and microeconometric methods.
\end{itemize}

\textbf{The Lucas critique, step by step.}
\st{1} A reduced-form regression of output on taxes is estimated from past data,
$Y_t=\alpha+\delta T_t+\varepsilon_t$, and used for a policy counterfactual.
\st{2} But taxes are set by a policy rule, $T_t=\tau Y_t$, and agents optimize
knowing the rule. So the coefficient is itself a function of the rule,
\[
Y_t=\alpha+\delta(\tau)\,T_t(\tau)+\varepsilon_t .
\]
\st{3} Changing the policy from $\tau$ to $\tau'$ changes $\delta(\tau)$ to
$\delta(\tau')$. A counterfactual that keeps the old $\hat\delta$ is
conceptually wrong.
\st{4} The remedy is to model the optimization problems (preferences,
technology, constraints) whose parameters do not move with policy. This is why
the rest of the course starts from household and planner problems.

% ======================================================================
\section*{Part 2. National accounts (Discussion 1)}
\addcontentsline{toc}{section}{Part 2. National accounts}

\textbf{Definition.} GDP is the market value of all \emph{final} goods and
services produced within a country in a given period. Factors of production are
long-lived inputs (labor, mines, machines). Intermediate goods are used up in
later stages (ore, steel). Final goods are not used up in later production
within the period (autos, new machines).

\textbf{Three equivalent approaches.}
\begin{enumerate}
\item \emph{Production.} $\text{GDP}=\sum_{\text{stages}}\big[\text{value of the stage's output}-\text{value of intermediates used}\big]$. Output includes unsold output (inventory). Adding gross output at every stage would double count intermediates.
\item \emph{Expenditure.} $Y_t=C_t+I_t+G_t+NX_t$. $C$ is household spending (durables, nondurables, services, about two-thirds of US GDP). $I$ is fixed investment (non-residential structures, equipment, intellectual property, plus residential housing) plus inventory investment. $G$ is government \emph{purchases} only (transfers and interest on debt are excluded). $NX$ is exports minus imports.
\item \emph{Income.} GDP is the sum of payments to factors plus indirect taxes. National income (compensation, proprietors' income, rental income, corporate profits, net interest) relates to GDP by $NI_t=GDP_t+NFP_t-D_t-NIT_t$.
\end{enumerate}

\textbf{Identities.}
\begin{align*}
\text{(A) expenditure}\quad & Y_t=C_t+I_t+G_t+NX_t\\
\text{(B) household}\quad & Y_t+TR_t+NFP_t+NTR_t=C_t+T_t+S_t\\
\text{(C) government}\quad & T_t+\dot M_t+\dot B^G_t=G_t+TR_t\\
\text{(D) balance of payments}\quad & X_t+NFP_t+NTR_t+CF_t=IM_t
\end{align*}

\textbf{Derivation of the saving--investment identity (E).}
\st{1} Start from (B) and replace $Y_t$ using (A):
\[
C_t+T_t+S_t=C_t+I_t+G_t+X_t-IM_t+TR_t+NFP_t+NTR_t .
\]
\st{2} Cancel $C_t$ on both sides:
\[
T_t+S_t=I_t+G_t+X_t-IM_t+TR_t+NFP_t+NTR_t .
\]
\st{3} From (D), $X_t+NFP_t+NTR_t-IM_t=-CF_t$. Substitute:
\[
T_t+S_t=I_t+G_t-CF_t+TR_t .
\]
\st{4} Solve for investment:
\ans{I_t=S_t+\underbrace{(T_t-G_t-TR_t)}_{\text{government surplus}}+CF_t}
Investment is financed by private saving, public saving, or foreign capital. In
a closed economy without government, $T=G=TR=CF=0$ and $Y_t=C_t+S_t$, so
$S_t=I_t$.

\textbf{CPI.} Inflation $\pi_{t+1}=(P_{t+1}-P_t)/P_t$ with $P_t$ the cost of a
fixed basket relative to the base year. Two biases make the CPI overstate
cost-of-living growth. Quality bias counts quality improvements as price
increases. Substitution bias ignores that consumers switch away from goods
that became relatively expensive.

\prob{1 (three approaches).} A mine produces ore worth 100, paying wages 60.
A mill turns all the ore into steel worth 250, paying wages 100. A factory turns
all the steel into autos worth 500, paying wages 150. Households buy autos
worth 400, the rest stay in the factory's inventory. Compute GDP all three ways.
\sol
\st{1 (production)} Value added at each stage is output minus intermediates.
\[
\underbrace{100-0}_{\text{mine}}+\underbrace{250-100}_{\text{mill}}+\underbrace{500-250}_{\text{factory}}=100+150+250=500 .
\]
\st{2 (expenditure)} Final goods are autos. Sold to households is $C=400$, unsold
is inventory investment $I=100$.
\[
C+I=400+100=500 .
\]
\st{3 (income)} Profits are value added minus wages, $100-60=40$,
$150-100=50$, $250-150=100$. Payments to factors:
\[
\underbrace{60+100+150}_{\text{wages }310}+\underbrace{40+50+100}_{\text{profits }190}=500 .
\]
\ans{\text{GDP}=500\text{ under all three approaches}}

\prob{2 (classification).} Say what happens to US and French GDP components.
(a) A French tourist buys and eats a hamburger in Ann Arbor. (b) A French firm
buys a US-made machine and uses it in the US. (c) A car made in the US in
November 2008 is sold to a US household in February 2009.
\sol
(a) US exports rise, so US $NX$ and $Y$ rise. French $C$ rises and French
imports rise by the same amount, so French $Y$ is unchanged.
(b) The machine is produced and installed in the US, so US $I$ and $Y$ rise.
French accounts do not change because the investment is not on French soil.
(c) In 2008 US $Y$ and $I$ (inventory) rise. In 2009 $C$ rises and inventory
investment falls by the same amount, so 2009 $Y$ is unchanged.

% ======================================================================
\section*{Part 3. The endowment economy (Lecture 2, Discussion 2, PSet 2)}
\addcontentsline{toc}{section}{Part 3. The endowment economy}

\subsection*{3.1 Setup}
\begin{itemize}
\item Discrete time $t=0,1,\dots,T$ with $T$ odd. One perishable good.
\item Two household types $e\in\{A,B\}$ with masses $n^A$, $n^B$. Endowments alternate:
\[
\begin{array}{c|cccccc}
 & 0&1&2&3&\cdots&T\\\hline
y^A_t&1&0&1&0&\cdots&0\\
y^B_t&0&1&0&1&\cdots&1
\end{array}
\]
\item Preferences $\sum_{t=0}^T\beta^t\ln c^e_t$ with $\beta\in(0,1)$, $\beta=1/(1+\rho)$, $\rho$ the rate of time preference.
\item Households can save ($s^e_t>0$) or borrow ($s^e_t<0$) in nominal bonds at nominal rate $i_t$. Goods cost $p_t$.
\end{itemize}
PSet 2 writes the agent index as $\ell$ instead of $e$. Nothing else changes.

\subsection*{3.2 From the nominal to the real budget constraint}
\st{1} Nominal constraints (dollars):
\[
p_0c^e_0+s^e_0\le p_0y^e_0,\qquad p_tc^e_t+s^e_t\le p_ty^e_t+(1+i_{t-1})s^e_{t-1}\quad(t>0).
\]
\st{2} Divide period $t$ by $p_t$ (dollars over dollars per basket gives baskets)
and define real bonds $b^e_t:=s^e_t/p_t$:
\[
c^e_t+b^e_t\le y^e_t+(1+i_{t-1})\frac{s^e_{t-1}}{p_t}
=y^e_t+(1+i_{t-1})\frac{p_{t-1}}{p_t}\,\frac{s^e_{t-1}}{p_{t-1}}
=y^e_t+(1+i_{t-1})\frac{p_{t-1}}{p_t}\,b^e_{t-1}.
\]
\st{3} Define inflation $1+\pi_t=p_t/p_{t-1}$ and the real rate
\[
1+r_{t-1}:=\frac{1+i_{t-1}}{1+\pi_t}\qquad\Longrightarrow\qquad i_t\approx r_t+\pi_{t+1}\quad\text{(Fisher equation).}
\]
\st{4} The real budget constraint is
\ans{c^e_t+b^e_t\le y^e_t+(1+r_{t-1})\,b^e_{t-1},\qquad b^e_{-1}=0.}
Only the real rate matters, because the household cares only about baskets.

\subsection*{3.3 Why the budget constraint binds}
Suppose at the optimum the constraint is slack at some date $\tilde t$. Then the
household can raise $c_{\tilde t}$ to
$c^{**}_{\tilde t}=y_{\tilde t}+(1+r_{\tilde t-1})b^*_{\tilde t-1}-b^*_{\tilde t}>c^*_{\tilde t}$
keeping all other choices fixed. Since $u'>0$, utility rises by
$\beta^{\tilde t}[u(c^{**}_{\tilde t})-u(c^*_{\tilde t})]>0$, contradicting
optimality. So with strictly increasing utility every budget constraint holds with equality.

\subsection*{3.4 No-Ponzi-game condition (finite horizon)}
$b^e_T\ge0$. The household cannot die in debt. Without it, every household
would borrow without limit in the last period.

\subsection*{3.5 Household problem and equilibrium definition}
Taking $\{r_t\}_{t=0}^{T-1}$ as given, household $e$ solves
\[
\max_{\{c^e_t,b^e_t\}_{t=0}^T}\sum_{t=0}^T\beta^t\ln c^e_t
\quad\text{s.t.}\quad c^e_t+b^e_t\le y^e_t+(1+r_{t-1})b^e_{t-1},\quad b^e_{-1}=0,\quad b^e_T\ge0 .
\]

\textbf{Definition (competitive equilibrium).} Prices $\{r_t\}_{t=0}^{T-1}$ and
allocations $\{c^A_t,c^B_t,b^A_t,b^B_t\}_{t=0}^T$ such that
\begin{enumerate}
\item taking $\{r_t\}$ as given, the allocations solve each household's problem,
\item the goods market clears, $n^Ac^A_t+n^Bc^B_t=n^Ay^A_t+n^By^B_t$ for all $t$,
\item the bond market clears, $n^Ab^A_t+n^Bb^B_t=0$ for all $t$.
\end{enumerate}
Write all three parts on the exam. Missing ``taking prices as given'' or one of
the market-clearing conditions is the usual way to lose points.

\subsection*{3.6 Solving the household problem}
\st{1 (Lagrangian)} Current-value form, multiplier $\beta^t\lambda^e_t$ on each
budget constraint and $\mu^e_T$ on $b^e_T\ge0$:
\[
\mathcal L^e=\sum_{t=0}^T\beta^t\Big[\ln c^e_t+\lambda^e_t\big(y^e_t+(1+r_{t-1})b^e_{t-1}-c^e_t-b^e_t\big)\Big]+\mu^e_Tb^e_T .
\]
The present-value multiplier is $\lambda^{e,PV}_t=\beta^t\lambda^e_t$.

\st{2 (FOC for $c^e_t$)} $c_t$ appears only in the date-$t$ term:
\[
\frac{\partial\mathcal L^e}{\partial c^e_t}=\beta^t\Big(\frac{1}{c^e_t}-\lambda^e_t\Big)=0
\quad\Longrightarrow\quad\frac{1}{c^e_t}=\lambda^e_t .\tag{1}
\]
\st{3 (FOC for $b^e_t$, $t<T$)} $b_t$ appears in the date-$t$ constraint with
coefficient $-1$ and in the date-$(t+1)$ constraint with coefficient $(1+r_t)$:
\[
\frac{\partial\mathcal L^e}{\partial b^e_t}=-\beta^t\lambda^e_t+\beta^{t+1}\lambda^e_{t+1}(1+r_t)=0
\quad\Longrightarrow\quad\lambda^e_t=\beta(1+r_t)\lambda^e_{t+1}.\tag{2}
\]
\st{4 (FOC for $b^e_T$)} $b_T$ appears in the last budget and in the no-Ponzi term:
\[
-\beta^T\lambda^e_T+\mu^e_T=0\quad\Longrightarrow\quad\mu^e_T=\beta^T\lambda^e_T .\tag{3}
\]
\st{5 (remaining KKT conditions)} Constraints (4) $y^e_t+(1+r_{t-1})b^e_{t-1}-c^e_t-b^e_t\ge0$
and (5) $b^e_T\ge0$. Slackness (6) $\lambda^e_t(\cdots)=0$ and (7) $\mu^e_Tb^e_T=0$.
Signs (8) $\lambda^e_t\ge0$ and (9) $\mu^e_T\ge0$.

\st{6 (implications)} Log utility forces $c^e_t>0$, so by (1) $\lambda^e_t=1/c^e_t>0$.
By (6) the budget binds:
\[
c^e_t+b^e_t=y^e_t+(1+r_{t-1})b^e_{t-1}.\tag{$*$}
\]
By (3), $\mu^e_T=\beta^T\lambda^e_T>0$, so by (7) $b^e_T=0$. This is the
terminal (transversality) condition. Nobody leaves unspent savings either.

\st{7 (Euler equation)} Substitute (1) into (2):
\ans{\frac{1}{c^e_t}=\beta(1+r_t)\frac{1}{c^e_{t+1}}\qquad\Longleftrightarrow\qquad c^e_{t+1}=\beta(1+r_t)c^e_t .\tag{$**$}}
The household problem is now two equations, $(*)$ and $(**)$, with two boundary
conditions, $b^e_{-1}=0$ and $b^e_T=0$.

\subsection*{3.7 The intertemporal budget constraint}
\st{1} Define the price of date-$t$ goods in units of date-0 goods,
$q_t:=1/\prod_{\tau=0}^{t-1}(1+r_\tau)$, with $q_0=1$. Note $q_{t+1}=q_t/(1+r_t)$.
\st{2} Multiply $(*)$ at date $t$ by $q_t$:
\[
q_tc_t+q_tb_t=q_ty_t+q_t(1+r_{t-1})b_{t-1}=q_ty_t+q_{t-1}b_{t-1}.
\]
\st{3} Sum from $t=0$ to $T$. The bond terms telescope:
\[
\sum_{t=0}^Tq_tc_t+\sum_{t=0}^Tq_tb_t=\sum_{t=0}^Tq_ty_t+\sum_{t=0}^{T}q_{t-1}b_{t-1}
\quad\Longrightarrow\quad
\sum_{t=0}^Tq_tc_t+q_Tb_T=\sum_{t=0}^Tq_ty_t+q_0(1+r_{-1})b_{-1}.
\]
\st{4} Use $b_{-1}=0$ and $b_T=0$:
\ans{\sum_{t=0}^T\frac{c^e_t}{\prod_{\tau=0}^{t-1}(1+r_\tau)}=\sum_{t=0}^T\frac{y^e_t}{\prod_{\tau=0}^{t-1}(1+r_\tau)}\qquad\text{lifetime consumption}=\text{lifetime income.}}

\subsection*{3.8 The equilibrium interest rate}
\st{1} Goods market clearing at $t$, then the Euler equation $c_t=c_{t+1}/[\beta(1+r_t)]$ for each type:
\begin{align*}
n^Ay^A_t+n^By^B_t&=n^Ac^A_t+n^Bc^B_t
=n^A\frac{c^A_{t+1}}{\beta(1+r_t)}+n^B\frac{c^B_{t+1}}{\beta(1+r_t)}\\
&=\frac{1}{\beta(1+r_t)}\big(n^Ac^A_{t+1}+n^Bc^B_{t+1}\big)
=\frac{1}{\beta(1+r_t)}\big(n^Ay^A_{t+1}+n^By^B_{t+1}\big).
\end{align*}
The last step uses market clearing at $t+1$.
\st{2} Solve for the rate:
\ans{1+r_t=\frac{1}{\beta}\,\frac{Y_{t+1}}{Y_t},\qquad Y_t:=n^Ay^A_t+n^By^B_t .}
This is the key general-equilibrium result. With log utility and a common
$\beta$, the rate depends only on aggregate endowment growth.
\st{3} In the alternating economy, $Y_t=n^A$ at even $t$ and $Y_t=n^B$ at odd $t$:
\[
1+r_t=\begin{cases}\dfrac1\beta\dfrac{n^B}{n^A}, & t\text{ even},\\[2ex]\dfrac1\beta\dfrac{n^A}{n^B}, & t\text{ odd}.\end{cases}
\]
The rate zigzags. When aggregate goods are scarce today relative to tomorrow,
the rate is high, because everybody would like to borrow against tomorrow and
the rate must rise until nobody does on net. With
$n^A=n^B$ the rate is constant at $1/\beta$.

\subsection*{3.9 Equilibrium allocations}
\st{1} Set $n:=n^B/n^A$. Then $q_0=1$, $q_1=\beta/n$, $q_2=q_1/(1+r_1)=(\beta/n)(\beta n)=\beta^2$, and in general
\[
q_t=\beta^t\ \text{ at even }t,\qquad q_t=\beta^t/n\ \text{ at odd }t .
\]
\st{2} The Euler equation gives $c_{t+1}=\beta(1+r_t)c_t$, so $c_{t+1}=nc_t$ after an even date and $c_{t+1}=c_t/n$ after an odd date. Consumption alternates between $c_{ev}$ and $c_{od}=n\,c_{ev}$.
\st{3} Value of the consumption stream:
\[
\sum_tq_tc_t=\sum_{t\text{ even}}\beta^tc_{ev}+\sum_{t\text{ odd}}\frac{\beta^t}{n}\,n c_{ev}=c_{ev}\sum_{t=0}^T\beta^t .
\]
\st{4} Value of endowments. Type $A$ gets 1 at even dates, type $B$ gets 1 at odd dates:
\[
\sum_tq_ty^A_t=\sum_{t\text{ even}}\beta^t,\qquad\sum_tq_ty^B_t=\frac1n\sum_{t\text{ odd}}\beta^t .
\]
\st{5} Use T2 with $T$ odd: $\sum_{t=0}^T\beta^t=\frac{1-\beta^{T+1}}{1-\beta}$, $\sum_{\text{even}}\beta^t=\frac{1-\beta^{T+1}}{1-\beta^2}$, $\sum_{\text{odd}}\beta^t=\beta\sum_{\text{even}}\beta^t$. Ratio
\[
\frac{\sum_{\text{even}}\beta^t}{\sum_{t}\beta^t}=\frac{1-\beta^{T+1}}{1-\beta^2}\cdot\frac{1-\beta}{1-\beta^{T+1}}=\frac{1-\beta}{(1-\beta)(1+\beta)}=\frac{1}{1+\beta}.
\]
\st{6} Set lifetime consumption equal to lifetime income for each type:
\ans{c^A_{ev}=\frac{1}{1+\beta},\quad c^A_{od}=\frac{n}{1+\beta},\qquad c^B_{ev}=\frac{\beta}{n(1+\beta)},\quad c^B_{od}=\frac{\beta}{1+\beta}.}
\st{7 (check)} Even dates: $n^A\frac{1}{1+\beta}+n^B\frac{\beta}{n(1+\beta)}=n^A\frac{1}{1+\beta}+n^A\frac{\beta}{1+\beta}=n^A=Y_{ev}$.
Odd dates: $n^A\frac{n}{1+\beta}+n^B\frac{\beta}{1+\beta}=\frac{n^B}{1+\beta}+\frac{\beta n^B}{1+\beta}=n^B=Y_{od}$. Both markets clear.
With $n^A=n^B$ ($n=1$) consumption is flat, $c^A=1/(1+\beta)$, $c^B=\beta/(1+\beta)$.
Type $A$ consumes more because it receives its endowment first.

\prob{3 (two-period economy, Discussion 2).} Two agents live for $t=1,2$ with
$\ln c_1+\beta\ln c_2$ and budgets $c_1+b=y_1$, $c_2=y_2+(1+r)b$. Write
$q:=1/(1+r)$. (a) Solve one household's problem. (b) Find the equilibrium $q$.
(c) Solve for $\beta=0.9$, $y^A=(3,1)$, $y^B=(1,1)$.
\sol
\st{1 (combine budgets)} From the second budget, $b=q(c_2-y_2)$. Substitute into
the first: $c_1+qc_2=y_1+qy_2$. The right side is wealth measured in date-1 goods.
\st{2 (Lagrangian)} $\mathcal L=\ln c_1+\beta\ln c_2+\lambda\big(y_1+qy_2-c_1-qc_2\big)$.
\[
\frac{1}{c_1}=\lambda,\qquad\frac{\beta}{c_2}=\lambda q\quad\Longrightarrow\quad\frac{1}{c_1}=\frac{\beta}{q}\frac{1}{c_2}=\beta(1+r)\frac{1}{c_2},\qquad c_2=\frac{\beta}{q}c_1 .
\]
\st{3 (plug into the budget)}
\[
c_1+q\frac{\beta}{q}c_1=y_1+qy_2\ \Longrightarrow\ (1+\beta)c_1=y_1+qy_2 .
\]
\ans{c_1=\frac{y_1+qy_2}{1+\beta},\qquad c_2=\frac{\beta(y_1+qy_2)}{q(1+\beta)},\qquad b=y_1-c_1=\frac{\beta y_1-qy_2}{1+\beta}.}
Check $b$: $y_1-\frac{y_1+qy_2}{1+\beta}=\frac{(1+\beta)y_1-y_1-qy_2}{1+\beta}=\frac{\beta y_1-qy_2}{1+\beta}$.
\st{4 (equilibrium price)} Bond market clearing $b^A+b^B=0$:
\[
\frac{\beta y^A_1-qy^A_2}{1+\beta}+\frac{\beta y^B_1-qy^B_2}{1+\beta}=0\ \Longrightarrow\ \beta Y_1-qY_2=0
\ \Longrightarrow\ \boxed{\ q=\beta\frac{Y_1}{Y_2},\qquad 1+r=\frac1\beta\frac{Y_2}{Y_1}\ }
\]
the two-period version of 3.8. With $Y_1=Y_2$, $1+r=1/\beta$.
\st{5 (numbers)} $Y_1=3+1=4$, $Y_2=1+1=2$, so $q=0.9\cdot4/2=1.8$ and $1+r=1/1.8=0.5556$. The real rate is negative because aggregate goods fall.
\begin{align*}
c^A_1&=\frac{3+1.8\cdot1}{1.9}=\frac{4.8}{1.9}=2.5263, & c^A_2&=\frac{0.9}{1.8}\,2.5263=1.2632, & b^A&=3-2.5263=0.4737,\\
c^B_1&=\frac{1+1.8\cdot1}{1.9}=\frac{2.8}{1.9}=1.4737, & c^B_2&=\frac{0.9}{1.8}\,1.4737=0.7368, & b^B&=1-1.4737=-0.4737.
\end{align*}
\st{6 (check)} $c^A_1+c^B_1=4=Y_1$, $c^A_2+c^B_2=2=Y_2$, $b^A+b^B=0$. Second budget for $A$:
$1+b^A/q=1+0.4737/1.8=1.2632$. Correct.

Discussion 2's special cases to remember. Endowments $(1,2)$ and $(2,1)$ give
$Y_1=Y_2$, $q=\beta$, and perfect smoothing for both. Identical endowments
$(1,1)$, $(1,1)$ give no trade. Endowments $(0,0)$ and $(2,2)$ also give no
trade because the rich agent gains nothing from trading with an agent who has
nothing to repay with.

\prob{4 (unequal masses, PSet 2).} In the alternating economy with $n^A=1$,
$n^B=2$, $\beta=0.9$, $T=7$, find $1+r_t$ and both consumption paths.
\sol
\st{1} $n=n^B/n^A=2$. Even $t$: $1+r_t=\frac{1}{0.9}\cdot2=2.2222$. Odd $t$: $1+r_t=\frac{1}{0.9}\cdot\frac12=0.5556$.
\st{2} From 3.9: $c^A_{ev}=\frac{1}{1.9}=0.5263$, $c^A_{od}=\frac{2}{1.9}=1.0526$, $c^B_{ev}=\frac{0.9}{2\cdot1.9}=0.2368$, $c^B_{od}=\frac{0.9}{1.9}=0.4737$.
\st{3 (check even date)} $1\cdot0.5263+2\cdot0.2368=1.0000=n^A$. Odd date: $1\cdot1.0526+2\cdot0.4737=2.0000=n^B$.

% ======================================================================
\section*{Part 4. Infinite horizon and uncertainty (Lecture 2, PSets 3 and 6)}
\addcontentsline{toc}{section}{Part 4. Infinite horizon and uncertainty}

\subsection*{4.1 Why we need a no-Ponzi condition}
With $T\to\infty$ there is no last period, so $b_T\ge0$ is gone. Consider
the scheme: borrow $x>0$ at $t=0$ and consume it, then each period roll over
the debt plus interest:
\[
b_0=-x,\quad b_1=-x(1+r_0),\quad b_2=-x(1+r_0)(1+r_1),\quad\dots,\quad b_T=-x\prod_{\tau=0}^{T-1}(1+r_\tau).
\]
Consumption after date 0 is unaffected, so this is a free lunch of $x$. It is
ruled out by the \textbf{no-Ponzi-game (NPG) condition}
\ans{\lim_{T\to\infty}\frac{b_T}{\prod_{\tau=0}^{T-1}(1+r_\tau)}\ge0 .}
In the scheme, $b_T/\prod(1+r_\tau)=-x<0$ for every $T$, so it is violated.

\subsection*{4.2 The transversality condition implies NPG with equality}
\st{1} The TVC is $\lim_{T\to\infty}\beta^T\lambda_Tb_T=0$. In words, at
infinity the present value of asset holdings is zero.
\st{2} Iterate the multiplier equation (2) backward, $\lambda_t=\lambda_{t-1}/[\beta(1+r_{t-1})]$:
\[
\lambda_t=\frac{\lambda_{t-1}}{\beta(1+r_{t-1})}=\frac{\lambda_{t-2}}{\beta^2(1+r_{t-1})(1+r_{t-2})}=\dots=\frac{\lambda_0}{\beta^t\prod_{\tau=0}^{t-1}(1+r_\tau)} .
\]
\st{3} Substitute into the TVC:
\[
0=\lim_{T\to\infty}\beta^T\frac{\lambda_0}{\beta^T\prod_{\tau=0}^{T-1}(1+r_\tau)}b_T=\lambda_0\lim_{T\to\infty}\frac{b_T}{\prod_{\tau=0}^{T-1}(1+r_\tau)} .
\]
\st{4} Divide by $\lambda_0=1/c_0>0$. The NPG condition holds with equality.
NPG is a constraint (no borrowing schemes). TVC is an optimality condition (no
wasted savings). Together they pin the limit to exactly zero.

\subsection*{4.3 Uncertainty and the stochastic Euler equation}
\st{1 (setup)} $y_t=\bar y+\varepsilon_t$ with $\varepsilon_t$ iid, mean zero. At date $t$ the household solves
\[
\max E_t\sum_{j=0}^\infty\beta^j\ln c_{t+j}\quad\text{s.t.}\quad c_{t+j}+b_{t+j}\le y_{t+j}+(1+r_{t+j-1})b_{t+j-1}\ \text{ and NPG.}
\]
\st{2 (Lagrangian)}
$\mathcal L=E_t\sum_{j=0}^\infty\beta^j\big[\ln c_{t+j}+\lambda_{t+j}\big(y_{t+j}+(1+r_{t+j-1})b_{t+j-1}-c_{t+j}-b_{t+j}\big)\big]$.
\st{3 (FOCs)} For $c_t$, known at $t$: $1/c_t=\lambda_t$. For $b_t$, which appears at $t$ and $t+1$:
\[
E_t\big[-\lambda_t+\beta(1+r_t)\lambda_{t+1}\big]=0\quad\Longrightarrow\quad\lambda_t=\beta(1+r_t)E_t[\lambda_{t+1}] .
\]
$r_t$ comes out of the expectation because the bond's return is known when it is bought (riskless bond).
\st{4} Combine:
\ans{\frac{1}{c_t}=\beta(1+r_t)\,E_t\Big[\frac{1}{c_{t+1}}\Big],\qquad E_t\Big[\frac{1}{c_{t+1}}\Big]=\int\frac{1}{c_{t+1}}g(\varepsilon_{t+1})\,d\varepsilon_{t+1}.}
General utility: $u'(c_t)=\beta(1+r_t)E_t[u'(c_{t+1})]$.

\subsection*{4.4 Pricing the bond when agents are identical (PSet 6)}
\st{1} Bonds are in zero net supply, $\sum b=0$. With identical agents every
agent holds the same amount, so $b_t=0$ for everyone.
\st{2} Then the budget gives $c_t=y_t$. Nobody trades. The rate is whatever
makes everybody content not to trade.
\st{3} Put $c_t=y_t$ into the Euler equation:
\ans{R_t=1+r_t=\frac{u'(y_t)}{\beta\,E_t[u'(y_{t+1})]} .}

\prob{5 (iid endowment).} Identical agents with log utility. Endowment is iid,
$y\in\{1,2\}$ with probability $\tfrac12$ each. Find $R$ in each state.
\sol
\st{1} $E[u'(y')]=\tfrac12\cdot\frac11+\tfrac12\cdot\frac12=\tfrac34$, the same in both states because $y$ is iid.
\st{2} $R(2)=\dfrac{1/2}{\beta\cdot3/4}=\dfrac{2}{3\beta}$ and $R(1)=\dfrac{1/1}{\beta\cdot3/4}=\dfrac{4}{3\beta}$.
\ans{R(y_h)=\frac{2}{3\beta}<\frac1\beta<R(y_l)=\frac{4}{3\beta}}
In a boom agents want to save for the expected fall, which pushes the rate down.
With $\beta=0.96$, $R(2)=0.6944$, $R(1)=1.3889$.

\prob{6 (PSet 6, Markov endowment).} Same, but $y\in\{y_h,y_l\}=\{2,1\}$
follows a Markov chain that stays in its state with probability $p$. Find $R_h$, $R_l$.
\sol
\st{1} In the high state: $E[u'(y')|y_h]=p\cdot\frac12+(1-p)\cdot1=\frac{p+2-2p}{2}=\frac{2-p}{2}$. So
\[
R_h=\frac{1/2}{\beta(2-p)/2}=\frac{1}{\beta(2-p)} .
\]
\st{2} In the low state: $E[u'(y')|y_l]=p\cdot1+(1-p)\cdot\frac12=\frac{2p+1-p}{2}=\frac{1+p}{2}$. So
\[
R_l=\frac{1}{\beta(1+p)/2}=\frac{2}{\beta(1+p)} .
\]
\ans{R_h=\frac{1}{\beta(2-p)},\qquad R_l=\frac{2}{\beta(1+p)}}
For $p<1$, $2-p>1$ and $(1+p)/2<1$, so $R_h<1/\beta<R_l$. At $p=1$ both equal
$1/\beta$ (no expected change). With $\beta=0.96$, $p=0.9$: $R_h=0.9470$, $R_l=1.0965$.
Recursive form (PSet 6): state $(b,y)$ with $b=R_{t-1}a_t$, Bellman
$V(b,y)=\max_{a'}\{u(b+y-a')+\beta E[V(R(y)a',y')|y]\}$. The current $y$ is both
income and the forecaster of $y'$, which is why it is a state.

% ======================================================================
\section*{Part 5. Log-linearization and time averaging (Discussion 3, PSet 3)}
\addcontentsline{toc}{section}{Part 5. Log-linearization and time averaging}

\subsection*{5.1 Notation}
$x_t$ is the series, $\bar x$ its trend or steady state, and
$\hat x_t:=\frac{x_t-\bar x}{\bar x}\approx\ln x_t-\ln\bar x$ the percent
deviation. So $x_t=\bar x(1+\hat x_t)\approx\bar xe^{\hat x_t}$.

\subsection*{5.2 Three methods on $Y_t=A_tK_t^\alpha N_t^{1-\alpha}$}
\textbf{Method 1 (substitute exponentials).}
\st{1} Replace each variable by $\bar xe^{\hat x_t}$:
$\bar Ye^{\hat Y_t}=\bar Ae^{\hat A_t}\,\bar K^\alpha e^{\alpha\hat K_t}\,\bar N^{1-\alpha}e^{(1-\alpha)\hat N_t}$.
\st{2} The steady state satisfies $\bar Y=\bar A\bar K^\alpha\bar N^{1-\alpha}$. Divide it out:
$e^{\hat Y_t}=e^{\hat A_t+\alpha\hat K_t+(1-\alpha)\hat N_t}$.
\st{3} Take logs: $\hat Y_t=\hat A_t+\alpha\hat K_t+(1-\alpha)\hat N_t$.

\textbf{Method 2 (first-order Taylor in levels).}
\st{1} $Y_t\approx\bar Y+\frac{\partial Y}{\partial A}(A_t-\bar A)+\frac{\partial Y}{\partial K}(K_t-\bar K)+\frac{\partial Y}{\partial N}(N_t-\bar N)$
with $\frac{\partial Y}{\partial A}=\frac{\bar Y}{\bar A}$, $\frac{\partial Y}{\partial K}=\alpha\frac{\bar Y}{\bar K}$, $\frac{\partial Y}{\partial N}=(1-\alpha)\frac{\bar Y}{\bar N}$ at the steady state.
\st{2} Subtract $\bar Y$ and divide by $\bar Y$:
$\frac{Y_t-\bar Y}{\bar Y}\approx\frac{A_t-\bar A}{\bar A}+\alpha\frac{K_t-\bar K}{\bar K}+(1-\alpha)\frac{N_t-\bar N}{\bar N}$, the same answer.

\textbf{Method 3 (log then total differential).}
\st{1} $\ln Y_t=\ln A_t+\alpha\ln K_t+(1-\alpha)\ln N_t$.
\st{2} Differentiate: $\frac{dY_t}{Y_t}=\frac{dA_t}{A_t}+\alpha\frac{dK_t}{K_t}+(1-\alpha)\frac{dN_t}{N_t}$.
\st{3} Evaluate at the steady state with $dx/\bar x=\hat x$. Same answer.
Method 3 is the fastest on an exam.

\subsection*{5.3 Six rules (PSet 3, Q3)}
\begin{align*}
&\text{(a) } y=cx: && \hat y=\hat x &&\text{exact}\\
&\text{(b) } y=x^a: && \hat y=a\hat x &&\text{(first order)}\\
&\text{(c) } y=xz: && \hat y=\hat x+\hat z\\
&\text{(d) } y=x/z: && \hat y=\hat x-\hat z\\
&\text{(e) } y=x+z: && \hat y=\tfrac{\bar x}{\bar y}\hat x+\tfrac{\bar z}{\bar y}\hat z &&\text{(weights are steady-state shares)}\\
&\text{(f) } y=x-z: && \hat y=\tfrac{\bar x}{\bar y}\hat x-\tfrac{\bar z}{\bar y}\hat z
\end{align*}
Derivation of (e) by Method 3: $dy=dx+dz$. Divide by $\bar y$ and write
$dx/\bar y=(\bar x/\bar y)(dx/\bar x)=(\bar x/\bar y)\hat x$.

\prob{7 (three model equations).} Log-linearize (a) $Y_t=C_t+I_t$,
(b) $K_{t+1}=(1-\delta)K_t+I_t$, (c) the CRRA Euler equation $C_t^{-\gamma}=\beta R\,C_{t+1}^{-\gamma}$ with $R$ constant.
\sol
(a) Rule (e): $\hat Y_t=\frac{\bar C}{\bar Y}\hat C_t+\frac{\bar I}{\bar Y}\hat I_t$.
With $\bar C/\bar Y=0.8$, $\hat C=0.01$, $\hat I=-0.02$, the approximation gives $0.8(0.01)+0.2(-0.02)=0.0040$. With hats as percent
deviations rule (e) is exact. Reading hats as log deviations
($C=0.8e^{0.01}$, $I=0.2e^{-0.02}$) the exact value is $\ln(1.00408)=0.00407$.

(b) Differentiate: $dK_{t+1}=(1-\delta)dK_t+dI_t$. Divide by $\bar K$:
$\hat K_{t+1}=(1-\delta)\hat K_t+\frac{\bar I}{\bar K}\hat I_t$. In steady state $\bar I=\delta\bar K$, so $\bar I/\bar K=\delta$ and
\ans{\hat K_{t+1}=(1-\delta)\hat K_t+\delta\hat I_t .}
(c) Take logs: $-\gamma\ln C_t=\ln\beta+\ln R-\gamma\ln C_{t+1}$. The steady
state has the same equation with bars, so subtract it:
$-\gamma\hat C_t=-\gamma\hat C_{t+1}$, that is $\hat C_{t+1}=\hat C_t$.
If $R_t$ varies, the same steps give $\hat C_{t+1}-\hat C_t=\frac1\gamma\hat R_t$, so the IES $1/\gamma$ scales the response.

\subsection*{5.4 Working's time-averaging problem}
Measured consumption is the average of two periods, observed every other period:
$x_{t+1}=\frac{c_{t+1}+c_t}{2}$, with true consumption a random walk $c_t=c_{t-1}+\varepsilon_t$.

\prob{8 (Working 1960).} Let $u_{t+1}=x_{t+1}-x_{t-1}$. Show $u$ is not white noise and find its first autocorrelation.
\sol
\st{1 (write $u$ in shocks)}
\begin{align*}
u_{t+1}&=\tfrac12\big[c_{t+1}+c_t-c_{t-1}-c_{t-2}\big]\\
&=\tfrac12\big[(c_t+\varepsilon_{t+1})+c_t-c_{t-1}-c_{t-2}\big]=\tfrac12\big[2c_t+\varepsilon_{t+1}-c_{t-1}-c_{t-2}\big]\\
&=\tfrac12\big[2(c_{t-1}+\varepsilon_t)+\varepsilon_{t+1}-c_{t-1}-c_{t-2}\big]=\tfrac12\big[c_{t-1}+2\varepsilon_t+\varepsilon_{t+1}-c_{t-2}\big]\\
&=\tfrac12\big[(c_{t-2}+\varepsilon_{t-1})+2\varepsilon_t+\varepsilon_{t+1}-c_{t-2}\big]=\tfrac12\big[\varepsilon_{t+1}+2\varepsilon_t+\varepsilon_{t-1}\big].
\end{align*}
\st{2 (variance)} Shocks are uncorrelated, so
$\Var(u_{t+1})=\frac14(1+4+1)\sigma^2=\frac32\sigma^2$.
\st{3 (covariance with the previous observation)} $u_{t-1}=\frac12[\varepsilon_{t-1}+2\varepsilon_{t-2}+\varepsilon_{t-3}]$. The only shared shock is $\varepsilon_{t-1}$, with weight $\frac12\cdot1$ in each:
\[
\Cov(u_{t+1},u_{t-1})=\tfrac12\cdot\tfrac12\,\E[\varepsilon_{t-1}^2]=\tfrac14\sigma^2 .
\]
\st{4 (autocorrelation)}
\ans{\rho_1=\frac{\sigma^2/4}{3\sigma^2/2}=\frac16\neq0}
So $x$ looks predictable even though $c$ is a random walk, and a naive test would wrongly reject the random walk.
\st{5 (the fix)} $u_{t-3}=\frac12[\varepsilon_{t-3}+2\varepsilon_{t-4}+\varepsilon_{t-5}]$ shares no shock with $u_{t+1}$, so $\Cov(u_{t+1},u_{t-3})=0$. Skipping one observation removes the spurious correlation. A simulation in the check script gives $\rho_1=0.166$ and $-0.0001$ for the skipped lag.

% ======================================================================
\section*{Part 6. Dynamic programming (Lecture 3, Discussion 4, PSets 3--4)}
\addcontentsline{toc}{section}{Part 6. Dynamic programming}

\subsection*{6.1 The two ways to write the problem}
\textbf{Problem 1.} $\max\sum_{t=0}^\infty\beta^t\hat F(y_t)$ s.t. $x_{t+1}=h(x_t,y_t)$, $x_{t+1}\in\Gamma(x_t)$,
with $x$ the state, $y$ the control, $h$ the transition, $\Gamma$ the feasible set.

\textbf{Problem 2.} Solve the transition for the control, $y_t=\hat h(x_t,x_{t+1})$,
and define $F(x_t,x_{t+1}):=\hat F(\hat h(x_t,x_{t+1}))$. Then
\[
\max_{\{x_{t+1}\in\Gamma(x_t)\}}\sum_{t=0}^\infty\beta^tF(x_t,x_{t+1}).
\]
The transition has been folded into the objective.

\textbf{Sequence Euler equation.} $x_{t+1}$ appears in two terms,
$\beta^tF(x_t,x_{t+1})$ and $\beta^{t+1}F(x_{t+1},x_{t+2})$. With an interior
solution, differentiate and divide by $\beta^t$:
\ans{F_2(x_t,x_{t+1})+\beta F_1(x_{t+1},x_{t+2})=0,\qquad\lim_{T\to\infty}\beta^TF_1(x_T,x_{T+1})\,x_T=0\ \text{(TVC).}}

\textbf{Example (growth planner).} $\max\sum\beta^tu(c_t)$ s.t.
$k_{t+1}=(1-\delta)k_t+i_t$, $c_t+i_t\le y_t$, $y_t\le f(k_t)$, $c_t\ge0$, $k_{t+1}\ge0$, $k_0$ given.
\st{1} Combine and use $u'>0$ to make it bind: $c_t+k_{t+1}=f(k_t)+(1-\delta)k_t$.
\st{2} Feasible set: $c_t\ge0$ gives the upper bound, so $\Gamma(k)=[0,\,f(k)+(1-\delta)k]$.
\st{3} $F(k,k')=u\big(f(k)+(1-\delta)k-k'\big)$, so
\[
F_1(k,k')=u'(c)\big(f'(k)+1-\delta\big),\qquad F_2(k,k')=-u'(c).
\]
\st{4} Plug into the Euler equation:
\[
-u'(c_t)+\beta u'(c_{t+1})\big(f'(k_{t+1})+1-\delta\big)=0\ \Longrightarrow\ \boxed{\ u'(c_t)=\beta u'(c_{t+1})\big(f'(k_{t+1})+1-\delta\big)\ }
\]

\subsection*{6.2 Deriving the Bellman equation}
\st{1} Define the value function $V(x_0):=\max_{\{x_{t+1}\in\Gamma(x_t)\}}\sum_{t=0}^\infty\beta^tF(x_t,x_{t+1})$.
\st{2} Split off the first term:
$V(x_0)=\max\big\{F(x_0,x_1)+\beta\sum_{t=1}^\infty\beta^{t-1}F(x_t,x_{t+1})\big\}$.
\st{3} Maximize in steps, first over $x_1$, then over everything after:
\[
V(x_0)=\max_{x_1\in\Gamma(x_0)}\Big\{F(x_0,x_1)+\beta\max_{\{x_{t+1}\in\Gamma(x_t)\}_{t\ge1}}\sum_{t=1}^\infty\beta^{t-1}F(x_t,x_{t+1})\Big\}.
\]
\st{4} Re-index the inner sum with $s=t-1$: it is $\sum_{s=0}^\infty\beta^sF(x_{s+1},x_{s+2})$, the same problem started from $x_1$. So it equals $V(x_1)$:
\[
V(x_0)=\max_{x_1\in\Gamma(x_0)}\big\{F(x_0,x_1)+\beta V(x_1)\big\}.
\]
\st{5} Drop time subscripts ($x$ today, $x'$ tomorrow):
\ans{V(x)=\max_{x'\in\Gamma(x)}\big\{F(x,x')+\beta V(x')\big\},\qquad g(x)=\arg\max_{x'\in\Gamma(x)}\big\{F(x,x')+\beta V(x')\big\}.}
This is a functional equation. The unknown is the function $V$. The policy
satisfies $V(x)=F(x,g(x))+\beta V(g(x))$. That the sequence problem and the
recursive problem give the same answer is the \textbf{principle of optimality}.
Growth model: $V(k)=\max_{k'\in[0,f(k)+(1-\delta)k]}\{u(f(k)+(1-\delta)k-k')+\beta V(k')\}$.

\subsection*{6.3 Properties of the value function (Lecture 3, p.~5)}
\begin{enumerate}
\item $V$ is often unique and can be approached by iterations (6.4).
\item If $F$ is strictly increasing in its first argument and $\Gamma$ is monotone ($x\le\tilde x\Rightarrow\Gamma(x)\subseteq\Gamma(\tilde x)$), then $V$ is strictly increasing.
\item If $F$ is jointly strictly concave and $\Gamma$ is convex ($x'\in\Gamma(x)$, $\tilde x'\in\Gamma(\tilde x)\Rightarrow\theta x'+(1-\theta)\tilde x'\in\Gamma(\theta x+(1-\theta)\tilde x)$ for $\theta\in(0,1)$), then $V$ is strictly concave and the policy is unique.
\item $V$ is sometimes differentiable, which gives the Euler equation (6.5). The policy is sometimes increasing.
\end{enumerate}

\subsection*{6.4 The Bellman operator is a contraction}
\textbf{Definitions.} $\Theta$ is the set of continuous real functions on a closed
bounded $X$. The Bellman operator is $(Tf)(x):=\max_{x'\in\Gamma(x)}\{F(x,x')+\beta f(x')\}$.
A metric $d$ on $M$ satisfies $d(x,y)=0\iff x=y$, symmetry, and the triangle
inequality. The sup-norm metric is $d(f,g)=\sup_x|f(x)-g(x)|$. $T$ is a
\textbf{contraction with modulus} $\beta\in[0,1)$ if $d(Tf,Tg)\le\beta\,d(f,g)$ for all $f,g$.

\textbf{Prop 1 (Banach).} If $(\Theta,d)$ is a non-empty complete metric space and $T$ a contraction with modulus $\beta$, then (i) $T$ has a unique fixed point $f^*=Tf^*$, and (ii) $d(Tf_0,f^*)\le\beta\,d(f_0,f^*)$ for any $f_0$.

\textbf{Prop 2 (Blackwell's sufficient conditions).} If $T$ satisfies
(i) \emph{monotonicity}, $f\le g$ everywhere $\Rightarrow Tf\le Tg$ everywhere, and
(ii) \emph{discounting}, there is $\beta\in[0,1)$ with $T(f+k)\le Tf+\beta k$ for every constant $k\ge0$,
then $T$ is a contraction with modulus $\beta$.

\textbf{Verifying Blackwell for the Bellman operator.}
\st{1 (monotonicity)} Take $f\le g$. For every $x'$, $F(x,x')+\beta f(x')\le F(x,x')+\beta g(x')$ because $\beta>0$. The max of a smaller function is smaller:
\[
Tf(x)=\max_{x'}\{F(x,x')+\beta f(x')\}\le\max_{x'}\{F(x,x')+\beta g(x')\}=Tg(x).
\]
\st{2 (discounting)} For a constant $k$,
\[
T(f+k)(x)=\max_{x'}\{F(x,x')+\beta f(x')+\beta k\}=\max_{x'}\{F(x,x')+\beta f(x')\}+\beta k=Tf(x)+\beta k .
\]
The constant comes out of the max because it does not depend on $x'$. This holds with equality, which needs $\beta<1$ to count.
\st{3} So $T$ is a contraction. Bounded continuous functions with the sup norm form a complete space, so Banach applies: $V$ exists, is unique, and
\[
d(T^nV_0,V)\le\beta\,d(T^{n-1}V_0,V)\le\dots\le\beta^nd(V_0,V)\to0 .
\]
\textbf{Algorithm (VFI).} Guess $V_0$ (for example $V_0\equiv0$), iterate $V_{n+1}=TV_n$ until $\sup|V_{n+1}-V_n|<$ tol.
\textbf{Caveat.} If $F$ is unbounded (log utility, $f(k)=k^\alpha$ on $[0,\infty)$), $\Theta$ contains unbounded functions, completeness can fail, and Banach need not apply directly.
Chain to remember: Blackwell $\Rightarrow$ contraction $\Rightarrow$ Banach $\Rightarrow$ $V$ unique and VFI converges.

\subsection*{6.5 The Euler equation from the Bellman equation}
This is Discussion 4's four-step recipe: (1) state and control, (2) Bellman, (3) FOC, (4) envelope.
\st{1 (FOC)} Differentiate the right side of the Bellman equation in $x'$:
$F_2(x,x')+\beta V'(x')=0$. The policy solves it: $F_2(x,g(x))+\beta V'(g(x))=0$.
\st{2 (envelope)} Differentiate $V(x)=F(x,g(x))+\beta V(g(x))$ in $x$ with the chain rule:
\begin{align*}
V'(x)&=F_1(x,g(x))+F_2(x,g(x))g'(x)+\beta V'(g(x))g'(x)\\
&=F_1(x,g(x))+\underbrace{\big[F_2(x,g(x))+\beta V'(g(x))\big]}_{=0\text{ by the FOC}}g'(x)=F_1(x,g(x)).
\end{align*}
\st{3 (shift forward one period)} $V'(x')=F_1(x',g(x'))=F_1(g(x),g(g(x)))$.
\st{4 (substitute into the FOC)}
\ans{F_2\big(x,g(x)\big)+\beta F_1\big(g(x),g(g(x))\big)=0\qquad\text{(functional Euler equation)}}
It is the sequence Euler equation with $x_{t+1}=g(x_t)$, $x_{t+2}=g(g(x_t))$.

\prob{9 (growth model Euler and steady state).} For
\[
V(k)=\max_{k'}\big\{u(f(k)+(1-\delta)k-k')+\beta V(k')\big\},
\]
derive the Euler equation by the four-step recipe and find the steady state for $f(k)=k^\alpha$.
\sol
\st{1} State $k$, control $k'$ (then $c$ follows from the resource constraint).
\st{2} FOC in $k'$: $-u'(c)+\beta V'(k')=0$, so $u'(c)=\beta V'(k')$.
\st{3} Envelope: $V'(k)=u'(c)\big(f'(k)+1-\delta\big)$. Shift forward: $V'(k')=u'(c')\big(f'(k')+1-\delta\big)$.
\st{4} Substitute: $u'(c)=\beta u'(c')\big(f'(k')+1-\delta\big)$.
\st{5 (steady state)} $c=c'$ cancels $u'$: $1=\beta(\alpha k^{\alpha-1}+1-\delta)$, so
$\alpha k^{\alpha-1}=\frac1\beta-1+\delta$ and
\ans{k^*=\Big(\frac{\alpha}{1/\beta-1+\delta}\Big)^{\frac{1}{1-\alpha}}}
With $\delta=1$: $k^*=(\alpha\beta)^{1/(1-\alpha)}$. For $\alpha=0.3$, $\beta=0.6$ this is $0.18^{1/0.7}=0.0863$.

\prob{10 (guess and verify, Discussion 4).} Log utility, $\delta=1$, $f(k)=k^\alpha$. Guess $V(k)=E+F\ln k$. Find $E$, $F$, and the policy.
\sol
\st{1} Plug the guess into the right side: $V(k)=\max_{k'}\{\ln(k^\alpha-k')+\beta(E+F\ln k')\}$.
\st{2 (FOC)}
\[
\frac{1}{k^\alpha-k'}=\frac{\beta F}{k'}\ \Longrightarrow\ k'=\beta F(k^\alpha-k')\ \Longrightarrow\ k'(1+\beta F)=\beta Fk^\alpha\ \Longrightarrow\ k'=\frac{\beta F}{1+\beta F}k^\alpha .
\]
Then $c=k^\alpha-k'=\frac{1}{1+\beta F}k^\alpha$.
\st{3 (substitute back)}
\begin{align*}
V(k)&=\ln\Big(\frac{k^\alpha}{1+\beta F}\Big)+\beta E+\beta F\ln\Big(\frac{\beta Fk^\alpha}{1+\beta F}\Big)\\
&=\alpha\ln k-\ln(1+\beta F)+\beta E+\beta F\ln(\beta F)-\beta F\ln(1+\beta F)+\alpha\beta F\ln k\\
&=\underbrace{\beta E+\beta F\ln(\beta F)-(1+\beta F)\ln(1+\beta F)}_{\text{must equal }E}+\underbrace{\alpha(1+\beta F)}_{\text{must equal }F}\ln k .
\end{align*}
\st{4 (match the $\ln k$ coefficient)} $F=\alpha+\alpha\beta F\Rightarrow F(1-\alpha\beta)=\alpha\Rightarrow F=\frac{\alpha}{1-\alpha\beta}$.
Then $\beta F=\frac{\alpha\beta}{1-\alpha\beta}$ and $1+\beta F=\frac{1-\alpha\beta+\alpha\beta}{1-\alpha\beta}=\frac{1}{1-\alpha\beta}$.
\st{5 (policy)} $\frac{\beta F}{1+\beta F}=\frac{\alpha\beta}{1-\alpha\beta}\cdot(1-\alpha\beta)=\alpha\beta$:
\ans{k'=\alpha\beta k^\alpha,\qquad c=(1-\alpha\beta)k^\alpha}
\st{6 (match the constant)} $E=\beta E+\beta F\ln(\beta F)-(1+\beta F)\ln(1+\beta F)$. Substitute
$\ln(\beta F)=\ln\frac{\alpha\beta}{1-\alpha\beta}=\ln(\alpha\beta)-\ln(1-\alpha\beta)$ and
$\ln(1+\beta F)=-\ln(1-\alpha\beta)$:
\begin{align*}
(1-\beta)E&=\frac{\alpha\beta}{1-\alpha\beta}\big[\ln(\alpha\beta)-\ln(1-\alpha\beta)\big]+\frac{1}{1-\alpha\beta}\ln(1-\alpha\beta)\\
&=\frac{\alpha\beta}{1-\alpha\beta}\ln(\alpha\beta)+\frac{1-\alpha\beta}{1-\alpha\beta}\ln(1-\alpha\beta)
=\frac{\alpha\beta}{1-\alpha\beta}\ln(\alpha\beta)+\ln(1-\alpha\beta).
\end{align*}
\ans{E=\frac{1}{1-\beta}\Big[\ln(1-\alpha\beta)+\frac{\alpha\beta}{1-\alpha\beta}\ln(\alpha\beta)\Big],\qquad F=\frac{\alpha}{1-\alpha\beta}}
With $\alpha=0.3$, $\beta=0.6$: $F=0.3659$, $E=-1.4372$. The check script confirms the Bellman residual is $3\times10^{-11}$.

\prob{11 (VFI by hand, analytically).} Same model, start from $V_0\equiv0$. Find $V_1$, $V_2$, and the general pattern.
\sol
\st{1} $V_1(k)=\max_{k'}\ln(k^\alpha-k')$. With no future value, save nothing: $k'=0$, $V_1(k)=\alpha\ln k$.
\st{2} Suppose $V_n(k)=a_n+b_n\ln k$. Then $V_{n+1}(k)=\max_{k'}\{\ln(k^\alpha-k')+\beta a_n+\beta b_n\ln k'\}$, and the FOC is Problem 10's with $F$ replaced by $b_n$:
\[
k'=\frac{\beta b_n}{1+\beta b_n}k^\alpha .
\]
\st{3} Substituting back, the $\ln k$ coefficient becomes $b_{n+1}=\alpha(1+\beta b_n)$ (same algebra as Step 3 of Problem 10). With $b_0=0$:
\[
b_1=\alpha,\quad b_2=\alpha(1+\alpha\beta),\quad b_3=\alpha(1+\alpha\beta+\alpha^2\beta^2),\quad b_n=\alpha\frac{1-(\alpha\beta)^n}{1-\alpha\beta}\to\frac{\alpha}{1-\alpha\beta}=F .
\]
\st{4 (numbers, $\alpha=0.3$, $\beta=0.6$)}
\[
\begin{array}{c|ccccc|c}
n&1&2&3&4&5&\infty\\\hline
b_n&0.3000&0.3540&0.3637&0.3655&0.3658&0.3659\\
\text{savings rate }\frac{\beta b_{n-1}}{1+\beta b_{n-1}}&0&0.1525&0.1752&0.1791&0.1799&0.18=\alpha\beta
\end{array}
\]
$V_n$ is the value with $n$ periods left. Each iteration adds one more period of future, so capital is worth more and the savings rate rises toward $\alpha\beta$.

\prob{12 (cake eating, PSet 3 Q1).} $\max\sum\beta^tu(c_t)$ with $u(c)=\frac{c^{1-\gamma}}{1-\gamma}$, $g_{t+1}=g_t-c_t$. Find the policy and value.
\sol
\st{1} Bellman: $V(g)=\max_{g'}\{u(g-g')+\beta V(g')\}$. FOC $u'(c)=\beta V'(g')$. Envelope $V'(g)=u'(c)$. Euler: $c_t^{-\gamma}=\beta c_{t+1}^{-\gamma}$, so $c_{t+1}=\beta^{1/\gamma}c_t=:\theta c_t$.
\st{2} Budget: total consumption equals the cake, $\sum_tc_t=g_0$, so $c_0\sum_t\theta^t=g_0$ and $c_0=(1-\theta)g_0$. The same holds at every date:
\ans{c=(1-\theta)g,\qquad g'=\theta g,\qquad\theta=\beta^{1/\gamma}}
\st{3 (value)} $V(g)=\sum_t\beta^t\frac{[(1-\theta)\theta^tg]^{1-\gamma}}{1-\gamma}=u(g)(1-\theta)^{1-\gamma}\sum_t(\beta\theta^{1-\gamma})^t$.
Now $\beta\theta^{1-\gamma}=\beta\cdot\beta^{(1-\gamma)/\gamma}=\beta^{1/\gamma}=\theta$, so the sum is $1/(1-\theta)$ and
\ans{V(g)=u(g)\,(1-\theta)^{-\gamma}}
With $\gamma=2$, $\beta=0.95$: $\theta=0.9747$, so 2.53\% of the cake is eaten each period.

\prob{13 (stochastic consumption-saving, Lecture 3).} $\max E_0\sum\beta^tu(c_t)$ s.t. $c_t+\frac{a_{t+1}}{1+r}=y_t+a_t$, with $y_t$ a two-state Markov chain. Derive the Euler equation from the Bellman equation.
\sol
\st{1} States $(a,y)$. $a$ is wealth and $y$ is income, which also forecasts $y'$. Control $a'$.
\st{2} Bellman:
\begin{align*}
V(a,y)&=\max_{a'}\Big\{u\Big(a+y-\frac{a'}{1+r}\Big)+\beta E[V(a',y')|y]\Big\},\\
E[V(a',y')|y]&=P(y_h|y)V(a',y_h)+P(y_l|y)V(a',y_l).
\end{align*}
\st{3 (FOC)} $0=-\frac{1}{1+r}u'(c)+\beta E\big[\frac{\partial V(a',y')}{\partial a'}\big|y\big]$.
\st{4 (envelope)} With $a'=g(a,y)$,
\[
\frac{\partial V}{\partial a}=u'(c)+\underbrace{\Big[-\frac{u'(c)}{1+r}+\beta E\Big[\frac{\partial V(a',y')}{\partial a'}\Big|y\Big]\Big]}_{=0\text{ by the FOC}}\frac{\partial g}{\partial a}=u'(c).
\]
\st{5} Shift forward, $\partial V(a',y')/\partial a'=u'(c')$, and substitute into the FOC:
\ans{u'(c)=\beta(1+r)\,E[u'(c')|y]=\beta(1+r)\sum_{y'\in\{y_l,y_h\}}P(y'|y)\,u'\big(h(a',y')\big)}
with $h(a,y):=a+y-g(a,y)/(1+r)$ the consumption policy.

% ======================================================================
\section*{Part 7. Numerical dynamic programming (Discussions 5--6, PSets 4--5)}
\addcontentsline{toc}{section}{Part 7. Numerical dynamic programming}

\subsection*{7.1 VFI on a grid}
Three numerical problems. \emph{Function approximation}: store $V$ on a grid,
an $n_a\times n_y$ matrix. \emph{Optimization}: grid search, pick the best
column. \emph{Integration}: with discrete shocks the expectation is a weighted sum.
\begin{enumerate}
\item Grids $A=\{a_1<\dots<a_{n_a}\}$ with $a_1=a_{\min}$ (this imposes the borrowing limit) and $Y=\{y_H,y_L\}$.
\item Guess $V_0(a,y)=0$.
\item Update $V_{j+1}(a,y)=\max_{a'\in A}\{u(y+a-\frac{a'}{1+r})+\beta[pV_j(a',y_H)+(1-p)V_j(a',y_L)]\}$. If consumption is negative, set utility to $-\infty$.
\item Stop when $\max|V_{j+1}-V_j|<$ tol. Read off $g_a$ as the argmax and $g_c$ from the budget.
\end{enumerate}
VFI is slow compared with Euler-equation methods but always converges under
Blackwell and works even when the Euler equation fails (corners, kinks,
discrete choices).

\prob{14 (two iterations by hand).} Growth model, $\ln c$, $\delta=1$, $k'=k^\alpha-c$, $\alpha=0.3$, $\beta=0.6$, grid $\{0.1,0.2,0.3\}$, $V_0\equiv0$.
\sol
\st{1 (output)} $0.1^{0.3}=0.5012$, $0.2^{0.3}=0.6170$, $0.3^{0.3}=0.6968$.
\st{2 (payoff matrix $\ln(k^\alpha-k')$, row $k$, column $k'$)}
\[
\begin{array}{c|ccc}
k\backslash k'&0.1&0.2&0.3\\\hline
0.1&\ln0.4012=-0.9133&\ln0.3012=-1.2000&\ln0.2012=-1.6035\\
0.2&\ln0.5170=-0.6596&\ln0.4170=-0.8746&\ln0.3170=-1.1487\\
0.3&\ln0.5968=-0.5161&\ln0.4968=-0.6995&\ln0.3968=-0.9242
\end{array}
\]
\st{3 (iteration 1)} $V_0=0$, so $V_1$ is the row max: $V_1=(-0.9133,-0.6596,-0.5161)$, policy $k'=0.1$ everywhere.
\st{4 (iteration 2)} Add $0.6V_1(k')=(-0.5480,-0.3958,-0.3097)$ to each row.
\[
\renewcommand{\arraystretch}{1.3}
\begin{array}{c|ccc|c}
k\backslash k'&0.1&0.2&0.3&V_2\\\hline
0.1&\begin{array}{c}-0.9133-0.5480\\=-1.4613\end{array}&\begin{array}{c}-1.2000-0.3958\\=-1.5958\end{array}&\begin{array}{c}-1.6035-0.3097\\=-1.9132\end{array}&-1.4613\\
0.2&\begin{array}{c}-0.6596-0.5480\\=-1.2076\end{array}&\begin{array}{c}-0.8746-0.3958\\=-1.2704\end{array}&\begin{array}{c}-1.1487-0.3097\\=-1.4584\end{array}&-1.2076\\
0.3&\begin{array}{c}-0.5161-0.5480\\=-1.0641\end{array}&\begin{array}{c}-0.6995-0.3958\\=-1.0953\end{array}&\begin{array}{c}-0.9242-0.3097\\=-1.2339\end{array}&-1.0641
\end{array}
\]
Policy still $k'=0.1$. This matches theory: the true policy $0.18k^{0.3}$ is at most about $0.125$ on this grid, closest to $0.1$.

\subsection*{7.2 Discretizing an AR(1) (Discussion 6)}
$z'=(1-\rho)\bar z+\rho z+\sigma\varepsilon'$, $\varepsilon'\sim N(0,1)$.
Stationary standard deviation $\sigma_z=\sigma/\sqrt{1-\rho^2}$.
Derivation: $\Var(z)=\rho^2\Var(z)+\sigma^2\Rightarrow\Var(z)=\sigma^2/(1-\rho^2)$.

\textbf{Tauchen (1986).}
\st{1 (grid)} $N$ equally spaced points on $[\bar z-m\sigma_z,\ \bar z+m\sigma_z]$, step $w=2m\sigma_z/(N-1)$.
\st{2 (probabilities)} Give point $z_\ell$ all the conditional mass within $w/2$ of it. Conditional on $z_j$, $z'$ is normal with mean $(1-\rho)\bar z+\rho z_j$ and sd $\sigma$:
\[
p_{j\ell}=\begin{cases}
\Phi\Big(\dfrac{z_1+w/2-(1-\rho)\bar z-\rho z_j}{\sigma}\Big), & \ell=1,\\[2ex]
\Phi\Big(\dfrac{z_\ell+w/2-(1-\rho)\bar z-\rho z_j}{\sigma}\Big)-\Phi\Big(\dfrac{z_\ell-w/2-(1-\rho)\bar z-\rho z_j}{\sigma}\Big), & 1<\ell<N,\\[2ex]
1-\Phi\Big(\dfrac{z_N-w/2-(1-\rho)\bar z-\rho z_j}{\sigma}\Big), & \ell=N.
\end{cases}
\]

\textbf{Rouwenhorst (1995).}
\st{1 (grid)} $N$ equally spaced points on $[\bar z-v,\bar z+v]$ with $v=\sigma_z\sqrt{N-1}$.
\st{2 (matrix)} Start from $P_2=\begin{pmatrix}p&1-p\\1-q&q\end{pmatrix}$ and build
\[
P_n=p\begin{pmatrix}P_{n-1}&\mathbf0\\\mathbf0'&0\end{pmatrix}+(1-p)\begin{pmatrix}\mathbf0&P_{n-1}\\0&\mathbf0'\end{pmatrix}+(1-q)\begin{pmatrix}\mathbf0'&0\\P_{n-1}&\mathbf0\end{pmatrix}+q\begin{pmatrix}0&\mathbf0'\\\mathbf0&P_{n-1}\end{pmatrix},
\]
then divide every row except the first and last by 2.
\st{3 (calibration)} Autocorrelation is $p+q-1=\rho$. With $p=q$, $p=\frac{1+\rho}{2}$.
Stationary variance is $v^2/(N-1)$, which equals $\sigma_z^2$ by the choice of $v$.
Rouwenhorst matches mean, variance and autocorrelation for any $N$ and is
preferred for persistent processes ($\rho>0.9$).
After discretizing $\tilde z=\log z$, use $z_j=\exp(\tilde z_j)$. The matrix is unchanged.

\prob{15 (Tauchen, $N=3$).} $\rho=0.9$, $\sigma=0.1$, $\bar z=0$, $m=1$. Use
$\Phi(0.9177)=0.8206$, $\Phi(1.1471)=0.8743$, $\Phi(3.2118)=0.9993$.
\sol
\st{1} $\sigma_z=0.1/\sqrt{1-0.81}=0.1/0.4359=0.2294$. Grid $\{-0.2294,0,0.2294\}$, $w=0.2294$, $w/2=0.1147$.
\st{2 (row 2, $z_j=0$, mean 0)} Cutoffs $\pm0.1147/0.1=\pm1.1471$:
$p_{21}=\Phi(-1.1471)=1-0.8743=0.1257$, $p_{23}=0.1257$, $p_{22}=1-2(0.1257)=0.7487$.
\st{3 (row 1, $z_j=-0.2294$, mean $-0.2065$)}
$p_{11}=\Phi\big(\frac{-0.2294+0.1147+0.2065}{0.1}\big)=\Phi(0.9177)=0.8206$,
$p_{13}=1-\Phi\big(\frac{0.2294-0.1147+0.2065}{0.1}\big)=1-\Phi(3.2118)=0.0007$,
$p_{12}=1-0.8206-0.0007=0.1787$.
\st{4} Row 3 mirrors row 1.
\ans{P=\begin{pmatrix}0.8206&0.1787&0.0007\\0.1257&0.7487&0.1257\\0.0007&0.1787&0.8206\end{pmatrix}}
With only 3 points and $m=1$ the chain's stationary sd is 0.175 against the true 0.229. Tauchen is poor with few points and high persistence, which is why Rouwenhorst exists.

\prob{16 (Rouwenhorst, $N=3$).} Same process. Build $P_3$ and check the moments.
\sol
\st{1} $p=q=\frac{1+0.9}{2}=0.95$, $v=\sqrt2\cdot0.2294=0.3244$, grid $\{-0.3244,0,0.3244\}$.
\st{2 (stack the four pieces)} First row: from the $p$ block $(p^2,\ p(1-p),\ 0)$ and from the $(1-p)$ block $(0,\ (1-p)p,\ (1-p)^2)$. Sum $(p^2,\ 2p(1-p),\ (1-p)^2)$.
Middle row, before halving: $p(1-q)+(1-q)p=2p(1-q)$, $pq+(1-p)(1-q)+(1-q)(1-p)+qp=2pq+2(1-p)(1-q)$, $(1-p)q+q(1-p)=2q(1-p)$. Halve.
Last row by symmetry: $((1-q)^2,\ 2q(1-q),\ q^2)$.
\[
P_3=\begin{pmatrix}p^2&2p(1-p)&(1-p)^2\\p(1-q)&pq+(1-p)(1-q)&q(1-p)\\(1-q)^2&2q(1-q)&q^2\end{pmatrix}
=\begin{pmatrix}0.9025&0.0950&0.0025\\0.0475&0.9050&0.0475\\0.0025&0.0950&0.9025\end{pmatrix}
\]
\st{3 (moments)} The stationary distribution is $(\tfrac14,\tfrac12,\tfrac14)$. Variance $\tfrac14v^2+\tfrac14v^2=\tfrac12v^2=\tfrac12\cdot2\sigma_z^2=\sigma_z^2=0.0526$. Exact.
Conditional mean from the top row ($z=-v$): $E[z'|z=-v]=-vp^2+v(1-p)^2=-v(p-(1-p))(p+(1-p))=-v(2p-1)=-0.9v=-\rho v$. By symmetry $E[z'|z=v]=\rho v$. Exact.

% ======================================================================
\section*{Part 8. Consumption (Lecture 4)}
\addcontentsline{toc}{section}{Part 8. Consumption}

\subsection*{8.1 The permanent income hypothesis (PIH) under certainty}
\st{1 (setup)} $\max\sum_t\beta^tu(c_t)$, $u'>0$, $u''<0$, subject to
$c_t+\frac{a_{t+1}}{1+r_t}=y_t+a_t$, $a_0$ given, and NPG
$\lim_{T\to\infty}a_{T+1}/\prod_{\tau=0}^{T}(1+r_\tau)\ge0$.
Here $a$ is wealth at the start of the period. This timing differs from Lecture 2's $b$.
\st{2 (Euler)} $u'(c_t)=\beta(1+r_t)u'(c_{t+1})$.
\st{3 (intertemporal budget)}
\[
\sum_{j=0}^\infty\frac{c_{t+j}}{\prod_{k=0}^{j-1}(1+r_{t+k})}=a_t+\sum_{j=0}^\infty\frac{y_{t+j}}{\prod_{k=0}^{j-1}(1+r_{t+k})}\qquad\text{lifetime consumption}=\text{lifetime wealth.}
\]
\st{4 (smoothing)} If $\beta(1+r_t)=1$ then $u'(c_t)=u'(c_{t+1})$, and since $u'$ is strictly decreasing, $c_t=c_{t+1}$.

\subsection*{8.2 CRRA and the intertemporal elasticity of substitution}
\st{1} With $u'(c)=c^{-1/\sigma}$ the Euler equation is $c_t^{-1/\sigma}=\beta(1+r_t)c_{t+1}^{-1/\sigma}$.
\st{2} Rearrange: $(c_{t+1}/c_t)^{1/\sigma}=\beta(1+r_t)$, so
\ans{\frac{c_{t+1}}{c_t}=\big[\beta(1+r_t)\big]^\sigma,\qquad\ln\frac{c_{t+1}}{c_t}=\sigma\ln(1+r_t)+\sigma\ln\beta .}
\st{3} So $\dfrac{\partial\ln(c_{t+1}/c_t)}{\partial\ln(1+r_t)}=\sigma$. That is the definition of the IES. Estimates put $\sigma\in[0.1,0.5]$.
\st{4} With $\beta=1/(1+\rho)$: consumption grows iff $\frac{1+r_t}{1+\rho}>1$, that is iff $r_t>\rho$.

\subsection*{8.3 The marginal propensity to consume}
\st{1} Constant $R=1+r$. Euler, applied $j$ times: $c_{t+j}=(R\beta)^\sigma c_{t+j-1}=\big((R\beta)^\sigma\big)^jc_t$.
\st{2} Plug into the budget:
\[
\sum_{j=0}^\infty R^{-j}\big((R\beta)^\sigma\big)^jc_t=c_t\sum_{j=0}^\infty\big(R^{-1}(R\beta)^\sigma\big)^j=\frac{c_t}{1-R^{-1}(R\beta)^\sigma}.
\]
\st{3} Set equal to wealth:
\ans{c_t=\big(1-R^{-1}(R\beta)^\sigma\big)\Big[a_t+\sum_{j=0}^\infty R^{-j}y_{t+j}\Big]}
\st{4 (transitory MPC)} $y_t$ enters with weight $R^0=1$:
\[
\text{mpc}=\frac{\partial c_t}{\partial y_t}=1-R^{-1}(R\beta)^\sigma\ \overset{R\beta=1}{=}\ 1-\frac{1}{1+r}=\frac{r}{1+r}\approx r .
\]
Small, around 0.03. A one-time windfall is spread over the whole future.
\st{5 (permanent MPC)} Write $y_{t+j}=y^P+y^c_{t+j}$. A rise in $y^P$ raises every term:
$\sum_jR^{-j}y^P=y^P\frac{1}{1-R^{-1}}=y^P\frac{1+r}{r}$. So with $R\beta=1$,
\[
\frac{\partial c_t}{\partial y^P}=\Big(1-\frac{1}{1+r}\Big)\frac{1+r}{r}=\frac{r}{1+r}\cdot\frac{1+r}{r}=1 .
\]
\ans{\text{mpc}\approx r\ \text{for a transitory change},\qquad\text{mpc}=1\ \text{for a permanent change}}
The MPC depends on the persistence of the income change (Friedman 1957).

\subsection*{8.4 Rational expectations errors and the random walk (Hall 1978)}
\textbf{RE error.} $\varepsilon_{t+1}:=s_{t+1}-E_t[s_{t+1}]$ with $E_t[\cdot]=E[\cdot\,|I_t]$. Three properties:
\begin{enumerate}
\item Mean zero: $E_t[\varepsilon_{t+1}]=E_t[s_{t+1}]-E_t[s_{t+1}]=0$.
\item Orthogonal to information: for $x_\tau\in I_t$, $E_t[x_\tau\varepsilon_{t+1}]=x_\tau E_t[\varepsilon_{t+1}]=0$.
\item Unpredictable by regression: the slope of $\varepsilon_{t+1}$ on $x_t$ is $\frac{\Cov(\varepsilon_{t+1},x_t)}{\Var(x_t)}=\frac{E[E_t[\varepsilon_{t+1}x_t]]}{\Var(x_t)}=0$ by the law of iterated expectations and property 2.
\end{enumerate}
\textbf{Random walk derivation.}
\st{1} Stochastic Euler: $u'(c_t)=\beta(1+r)E_t[u'(c_{t+1})]$.
\st{2} Write $u'(c_{t+1})=E_t[u'(c_{t+1})]+\varepsilon_{t+1}$. Then
$u'(c_t)=\beta(1+r)\big[u'(c_{t+1})-\varepsilon_{t+1}\big]$.
\st{3} With $\beta(1+r)\approx1$: $u'(c_{t+1})\approx u'(c_t)+\varepsilon_{t+1}$. Marginal utility is a random walk.
\st{4} Linearize $u'$ around steady-state $c$: $u'(c_t)\approx u'(c)+u''(c)c(\ln c_t-\ln c)$. Plug into both sides of Step 3, cancel $u'(c)$, and divide by $u''(c)c$:
\ans{\ln c_{t+1}=\ln c_t+\eta_{t+1},\qquad\eta_{t+1}=\frac{\varepsilon_{t+1}}{c\,u''(c)},\qquad E_t[\ln c_{t+1}]=\ln c_t .}
\st{5 (Hall's test)} Estimate $\ln c_{t+1}=\ln c_t+\gamma x_t+\eta_{t+1}$ with $x_t\in I_t$ (lagged income, stock prices). The PIH says $H_0:\gamma=0$.
\textbf{Evidence.} Shea (1995) finds that predictable union wage changes move
consumption, which contradicts the PIH. Johnson, Parker and Souleles (2006)
use the randomized timing of the 2001 tax rebates (\$300--600):
$c_{i,t}-c_{i,t-1}=\beta R_{i,t}+\text{controls}+u_{i,t}$. The PIH predicts
$\beta=0$ for an anticipated payment ($r/(1+r)\approx r$ even if unanticipated).
They estimate $\hat\beta\approx0.2$--$0.3$, which rejects the PIH. Also cited:
Hsieh (2003), Kueng (2016), Kaplan and Violante (2014).

\subsection*{8.5 Precautionary saving}
\st{1} With $\beta(1+r)=1$ the Euler equation is $u'(c_t)=E_t[u'(c_{t+1})]$.
\st{2} Expand $u'(c_{t+1})$ to second order around $c_t$:
$u'(c_{t+1})\approx u'(c_t)+u''(c_t)(c_{t+1}-c_t)+\frac12u'''(c_t)(c_{t+1}-c_t)^2$.
\st{3} Take $E_t$ and use Step 1 to cancel $u'(c_t)$:
$0=u''(c_t)E_t[c_{t+1}-c_t]+\frac12u'''(c_t)E_t[(c_{t+1}-c_t)^2]$.
\st{4} Solve and divide both sides by $c_t$ (on the right, write $1/c_t=c_t/c_t^2$):
\ans{E_t\Big[\frac{c_{t+1}-c_t}{c_t}\Big]=\underbrace{-\frac{u'''(c_t)c_t}{u''(c_t)}}_{\text{relative prudence}}\cdot\frac12\,E_t\Big[\Big(\frac{c_{t+1}-c_t}{c_t}\Big)^2\Big]}
Since $u''<0$, expected growth is positive iff $u'''>0$ (prudence, Kimball 1990).
Consumption is pushed into the future, so the household saves more today. This is precautionary saving.
\st{5 (which utilities)} Quadratic $u=b_1c-\frac12b_2c^2$ has $u'''=0$, so there is no precautionary saving (certainty equivalence). CRRA has $u'''=\frac1\sigma(\frac1\sigma+1)c^{-1/\sigma-2}>0$, relative prudence $1+\frac1\sigma$.
\st{6 (picture)} $u'$ convex ($u'''>0$) gives, by Jensen, $E_t[u'(c_{t+1})]>u'(E_t[c_{t+1}])$. Then $u'(c_t)=E_t[u'(c_{t+1})]>u'(E_t[c_{t+1}])$, and since $u'$ is decreasing, $c_t<E_t[c_{t+1}]$.

\subsection*{8.6 Borrowing constraints}
\textbf{Two kinds of limit.} $a_{t+1}\ge\underline a$ with $\underline a\le0$.
An \emph{ad hoc} limit is just imposed. The \emph{natural} limit is the most a
household can repay even if it consumes zero forever.
\st{1} Future budget from $t+1$ on with $c\ge0$:
$a_{t+1}+\sum_{j\ge0}\frac{y_{t+1+j}}{(1+r)^j}=\sum_{j\ge0}\frac{c_{t+1+j}}{(1+r)^j}\ge0$.
\st{2} This must hold for the worst income path, $y=y_{\min}$ forever:
\ans{a_{t+1}\ge-y_{\min}\sum_{j=0}^\infty\frac{1}{(1+r)^j}=-\frac{1+r}{r}\,y_{\min}}
\textbf{KKT with the constraint.}
\st{1} $\mathcal L=E_0\sum_t\beta^t\big[u(c_t)+\lambda_t\big(y_t+a_t-c_t-\frac{a_{t+1}}{1+r}\big)+\mu_t(a_{t+1}-\underline a)\big]$.
\st{2} FOC in $c_t$: $u'(c_t)=\lambda_t$. FOC in $a_{t+1}$ ($a_{t+1}$ appears at $t$ and as $a_{t+1}$ in the date-$(t+1)$ budget):
$-\frac{\lambda_t}{1+r}+\mu_t+\beta E_t[\lambda_{t+1}]=0$.
\st{3} Multiply by $(1+r)$ and substitute $\lambda=u'$:
\ans{u'(c_t)=\beta(1+r)E_t[u'(c_{t+1})]+(1+r)\mu_t,\qquad\mu_t\ge0,\quad\mu_t(a_{t+1}-\underline a)=0 .}
\st{4 (unconstrained)} $a_{t+1}>\underline a\Rightarrow\mu_t=0$, the usual Euler equation.
\st{5 (constrained)} $a_{t+1}=\underline a$, $\mu_t\ge0$, so $u'(c_t)\ge\beta(1+r)E_t[u'(c_{t+1})]$. Consumption is lower than the household would like. From the budget, $c_t=y_t+a_t-\frac{\underline a}{1+r}$, so
\ans{\text{mpc}=\frac{\partial c_t}{\partial y_t}=1\ \text{when the constraint binds}}
This high MPC for constrained households is the core of HANK models.
\st{6 (self-insurance)} With $\beta(1+r)=1$, $u'(c_t)\ge E_t[u'(c_{t+1})]$, so marginal utility is a supermartingale. Under conditions the supermartingale convergence theorem gives that $u'(c_t)$ converges almost surely. With income risk that never dies out, consumption cannot settle at a finite level, so the only limit consistent with the budget is $u'(c_t)\to0$, and $c$ and $a$ grow without bound. Using quadratic utility, $b_1-b_2c_t\ge E_t[b_1-b_2c_{t+1}]$, which means $c_t\le E_t[c_{t+1}]$. This isolates the constraint-driven saving motive from prudence, since quadratic utility has none.

\prob{17 (MPCs).} (a) $\beta(1+r)=1$, $r=0.04$. Find the exact transitory MPC.
(b) $\beta=0.95$, $r=0.04$, $\sigma=0.5$. Find the transitory MPC. (c) Find the MPC out of a permanent change in (a).
\sol
(a) $1-\frac{1}{1.04}=\frac{0.04}{1.04}=0.0385$.
(b) $R\beta=1.04\cdot0.95=0.988$, $(0.988)^{0.5}=0.99398$, divided by $1.04$ gives $0.95575$, so $\text{mpc}=1-0.95575=0.0442$.
(c) $\frac{0.04}{1.04}\cdot\frac{1.04}{0.04}=1$.

\prob{18 (natural borrowing limit).} $y_{\min}=0.5$, $r=0.04$.
\sol
$-\frac{1.04}{0.04}\cdot0.5=-26\cdot0.5=-13$. The household can never owe more than 13.

\prob{19 (size of the precautionary motive).} Log utility, $\beta(1+r)=1$,
consumption growth risk with standard deviation 5\%. Approximate expected growth.
\sol
Relative prudence for log is $1+\frac11=2$. So $E[\Delta c/c]\approx2\cdot\frac12\cdot(0.05)^2=0.0025$, or 0.25\% a year.
An exact two-point check ($\pm5\%$ around the mean) gives $0.249\%$.

\prob{20 (a binding borrowing constraint).} Two periods, $\beta=1$, $r=0$, log utility,
$y_1=1$, $y_2=3$, no borrowing ($a_2\ge0$). Find $c_1$, $c_2$, $\mu$, and the MPC out of $y_1$.
\sol
\st{1} Unconstrained, Euler with $\beta(1+r)=1$ gives $c_1=c_2=\frac{1+3}{2}=2$, which needs $a_2=y_1-c_1=-1<0$. Not allowed.
\st{2} So the constraint binds: $a_2=0$, $c_1=y_1=1$, $c_2=y_2=3$.
\st{3} Multiplier from $u'(c_1)=\beta(1+r)u'(c_2)+(1+r)\mu$: $1=\frac13+\mu$, so $\mu=\frac23>0$, consistent with KKT.
\st{4} $c_1=y_1$ one for one, so the MPC is 1. Without the constraint it would be $\frac12$.

\prob{21 (RE errors).} Under the PIH with quadratic utility and $\beta(1+r)=1$, show $E_t[c_{t+1}]=c_t$ and explain why regressing $\Delta c_{t+1}$ on lagged income should give zero.
\sol
\st{1} $u'(c)=b_1-b_2c$, so the Euler equation $b_1-b_2c_t=E_t[b_1-b_2c_{t+1}]$ gives $c_t=E_t[c_{t+1}]$.
\st{2} So $\Delta c_{t+1}=c_{t+1}-E_t[c_{t+1}]$ is an RE error.
\st{3} Lagged income is in $I_t$. By property 3 the population slope of an RE error on any variable in $I_t$ is zero. A nonzero estimate (excess sensitivity) rejects the PIH.

% ======================================================================
\section*{Sources}
\addcontentsline{toc}{section}{Sources}
\begin{itemize}
\item Lecture notes, ECO387C Fall 2026: Intro to Macro (Part 1), An endowment economy (Parts 3--4), Dynamic Programming (Part 6, Problem 13), Consumption (Part 8). Transcriptions in \texttt{lectures/}.
\item Discussions 1--6 (Kim and Yan): NIPA (Part 2), two-period equilibrium (Problem 3), log-linearization and Working (Part 5), DP recipe, guess and verify, VFI (Problems 9--11), VFI algorithm (7.1), Tauchen and Rouwenhorst (7.2).
\item Problem sets 2--6: alternating economy (Problem 4), cake eating (Problem 12), log-linearization rules (5.3), VFI by hand (Problem 14), Markov bond pricing (Problem 6).
\item Azzimonti, Krusell, McKay and Mukoyama, \emph{Macroeconomics} (2024 draft). Sec.~4.2--4.4 (sequential and recursive methods), Sec.~5.2--5.3 (Arrow-Debreu and sequential equilibrium in the endowment economy), Sec.~7.1 (Markov chains, AR processes), Sec.~7.6 (incomplete markets), Sec.~9.3 (PIH, borrowing constraints, precautionary saving).
\item Working, H. (1960), ``Note on the correlation of first differences of averages in a random chain,'' \emph{Econometrica} 28(4), 916--918.
\end{itemize}

\end{document}
